\hspace{3mm}

\centerline{\bf \huge 11 класс}
\centerline{\normalsize{\it Работа рассчитана на \bf{210} минут}}

\hspace{3mm}

\problem{11.1}Аюка выписал на доску многочлен $19x^2 - 2023x + 2005$, а Мингиян возвёл этот многочлен в $2023$ степень, раскрыл скобки, привёл подобные и посчитал сумму коэффициентов этого многочлена. Какое число у него получилось?

\hspace{3mm}

Ответ: 1.

Решение:\\Сумма коэффициентов любого многочлена от одной переменной $x$ равна значению многочлена при $x = 1.$ 
Имеем: $(19\cdot1^2 - 2023\cdot1 + 2005)^{2023} = 1^{2023} = 1.$

\hspace{3mm}

Критерии:\\
Любое правильное решение с верным обоснованием: \textbf{7 баллов.}\\
Неверное решение или идеи, не продвигающие ход решения: \textbf{0 баллов.}

\hspace{3mm}


\problem{11.2}Алтан и Бамба записали по одному целому числу на доску и рядом написали сумму этих чисел. Затем Алтан умножил своё число на $X$, а число Бамбы на $Y$, сложил полученные произведения и оказалось, что новая сумма делится на сумму чисел, записанную ребятами в самом начале на доске. Бамба решил поступить по-своему и умножил своё число на $X$, а число Алтана на $Y$ и сложил полученные произведения. Докажите, что новая сумма Бамбы также делится на изначальную сумму, записанную ребятами на доске, если известно, что $X$, $Y$ - целые числа.

\hspace{3mm}

Решение:\\
Пусть $A$ - число Алтана, а $B$ - число Бамбы. Следовательно, изначальная сумма $A + B$. Новая сумма, которая получилась у Алтана, равна: $AX + BY$, а у Бамбы - $BX + AY$. Рассмотрим сумму новых сумм мальчиков, мы имеем:
$(AX + BY) + (BX + AY) = A(X + Y) + B(X + Y) = (X + Y)(A + B)$. Очевидно, что данная сумма сумм делится на изначальную сумму $(A + B),$ и первое слагаемое $(AX + BY)$ делится на $(A + B)$ из условия, значит, и второе слагаемое $(BX + AY)$ делится на $(A + B)$. Значит, новая сумма Бамбы делится на изначальную сумму, записанную на доске.

\newpage

\hspace{3mm}

Критерии:\\
Приведено полное обоснованное решение: \textbf{7 баллов.}\\
Приведено верное в целом рассуждение, содержащее незначительные пробелы или неточности: \textbf{5 баллов.}\\
Идеи, продвигающие ход решения, но недостаточные для получения верного ответа: \textbf{2 балла.}\\
Рассмотрены только частные случаи или конкретные примеры: \textbf{0 баллов.}\\
Задача не решена или решена неверно: \textbf{0 баллов.}

\problem{11.3}Джангар каждый будний день недели побеждает некоторое количество противников, а в выходные дни он отдыхает. Известно, что в каждый последующий день недели он побеждал не больше противников, чем в предыдущий. И за всю третью неделю декабря, начиная с понедельника, он победил 6000 врагов. Какое наименьшее суммарное количество врагов он мог победить за первый, третий и пятый дни в третью неделю декабря?

\hspace{3mm}

Ответ: 3000 врагов.

Решение:\\
\textit{Оценка:} Докажем, что в первый, третий и пятый день меньше 3000 побеждённых врагов быть не могло. 
Действительно,  пусть в эти дни побеждено меньше 3000 врагов. Так как во второй и в четвёртый день богатырь победил не больше, чем в первый и в третий, то во второй и в четвёртый также побеждено меньше 3000 врагов. Тем самым в сумме за пять дней побеждено меньше 6000 врагов, что противоречит условию. 

\textit{Пример:} Если в каждый из первых четырёх дней богатырь побеждал по 1500 врагов, а в пятый - никого, то условия задачи выполнены. И за первый, третий и пятый дни он победил ровно 3000 врагов.

\hspace{3mm}

Критерии:\\
Приведена и доказана верная оценка на 3000 врагов и приведён любой правильный пример на 3000 врагов: \textbf{7 баллов.}\\
Приведена и доказана верная оценка на 3000 врагов, но пример не приведён или приведён неправильный: \textbf{5 баллов.}\\
Приведён любой правильный пример - 3000 врагов, но не доказано, почему не может быть меньше: \textbf{2 балла.}

\newpage

\hspace{3mm}

\noindentПриведена неверная оценка или приведён неверный пример: \textbf{0 баллов.}
Верный или неверный ответ без обоснования: \textbf{0 баллов.} 

\hspace{3mm}

\problem{11.4}Известно что $a, b, c \geq 0$
Докажите, что $(8a + b + c)(a + \dfrac{1}{8}b + c)(a + b + \dfrac{1}{27} c ) \geq 9abc$

\hspace{3mm}

Решение:

\textit{Полезный факт: Неравенство о среднем (для трех переменных).} Пусть $x, y, z \geq 0.$ Тогда справедливо неравенство:
$x+y+z \geq 3\sqrt[3]{xyz}$

Применим трижды неравенство между средним арифметическим и средним геометрическим, нужно написать это неравенство.

\noindent$1) 8a + b + c \geq 3\sqrt[3]{8abc} = 6 \sqrt[3]{abc}$\\
$2) a + 1/8b + c \geq 3\sqrt[3]{a \dfrac{1}{8} b  c} = 3/2 \sqrt[3]{abc}$\\
$3) a + b + \dfrac{1}{27}c \geq 3\sqrt[3]{ab \dfrac{1}{27}c} = \dfrac{3}{3}\sqrt[3]{abc}$\\

Перемножим все три неравенства и получим требуемое.

\hspace{3mm}

Критерии:\\
Любое правильное решение с верным обоснованием: \textbf{7 баллов.}\\
Идеи, продвигающие ход решения, но недостаточные для получения верного ответа: \textbf{2 балла.}\\
Неверное решение или идеи, не продвигающие ход решения: \textbf{0 баллов.}

\hspace{3mm}

\problem{11.5}Дан остроугольный треугольник ABC. Окружности с центрами $A$ и $C$ проходят через точку
$B$, вторично пересекаются в точке $F$ и пересекают описанную около треугольника ABC окружность $\omega$
в точках $D$ и $E$. Отрезок $BF$ пересекает окружность $\omega$ в точке $O$. Докажите, что $O$ — центр описанной
окружности треугольника $DEF$.

\hspace{3mm}

Решение:\\
Докажем сначала, что $OE = OF$. Для этого не нужна окружность с центром в точке $A$. Рассмотрим чертеж без нее (см. рис). Пусть в окружности с центром C центральный угол $BCE$ равен $2\alpha$. Тогда вписанный угол $EFB$ равен $\alpha$. В окружности $\omega$ углы $BCE$ и $BOE$ - вписанные и опираются на одну дугу, значит, $\angle BOE = \angle BCE = 2\alpha$.

\newpage

\hspace{3mm}

\noindentУгол $BOE$ - внешний угол треугольника $EOF$. Следовательно, $\angle OEF = \angle BOE - \angle OFE = \alpha,$ то есть треугольник $EOF$ - равнобедренный: $OE = OF$. Аналогично, рассмотрев окружность с центром $A$, доказывается, что $OF = OD$. Следовательно, $O$ - центр описанной окружности треугольника $DEF$ 

\begin{figure}[h] 
    \includegraphics[scale=0.7]{Городская Олимпиада Элиста 2023/Решения/picture_11.5_sol.jpg}
\end{figure}


\hspace{3mm}

Критерии:\\
Приведено полное обоснованное решение: \textbf{7 баллов.}\\
Приведено верное в целом рассуждение, в котором допущены несущественные пробелы или неточности: \textbf{5 баллов.}\\
В работе есть верная идея решения, но до конца оно не доведено: \textbf{3 балла.}\\
Неверное решение или идеи, не продвигающие ход решения: \textbf{0 баллов.}

